\documentclass[11pt,a4paper]{article}
\usepackage[utf8]{inputenc}
\usepackage[T1]{fontenc}
\usepackage{lmodern}
\usepackage[spanish,es-nodecimaldot]{babel}
\usepackage{amsmath,amssymb}
\usepackage{booktabs,array,tabularx}
\usepackage[margin=2.3cm]{geometry}
\usepackage{enumitem}
\usepackage[most]{tcolorbox}
\usepackage[colorlinks=true,linkcolor=black]{hyperref}
\setlength{\parskip}{5pt}\setlength{\parindent}{0pt}
\setlist{itemsep=2pt,topsep=3pt}
\renewcommand{\arraystretch}{1.15}

\newcommand{\EstFcv}{4{,}13}
\newcommand{\EstFt}{1{,}37}
\newcommand{\EstFv}{43{,}5}
\newcommand{\EstPcv}{\ensuremath{<0{,}001}}
\newcommand{\EstPt}{0{,}182}
\newcommand{\EstRv}{0{,}50}
\newcommand{\LamInicios}{0{,}22}
\newcommand{\LamInicioslo}{0{,}00}
\newcommand{\LamIniciosup}{0{,}44}
\newcommand{\LamTasa}{0{,}31}
\newcommand{\LamTasalo}{\ensuremath{-}0{,}05}
\newcommand{\LamTasaup}{0{,}67}
\newcommand{\PerfilIniciosLog}{\ensuremath{-}20{,}4}
\newcommand{\PerfilIniciosUno}{\ensuremath{-}22{,}6}
\newcommand{\PerfilTasaLog}{\ensuremath{-}8{,}9}
\newcommand{\PerfilTasaUno}{\ensuremath{-}19{,}7}
\newcommand{\VolNormal}{0{,}23}
\newcommand{\VolRatio}{2{,}6}
\newcommand{\VolVolcker}{0{,}58}

\newcommand{\AdfTasaDif}{\ensuremath{-}5{,}66}
\newcommand{\AdfTasaNivTM}{\ensuremath{-}1{,}98}
\newcommand{\AdfTasaNivTMp}{0{,}298}
\newcommand{\AdfTasaNivTT}{\ensuremath{-}1{,}96}
\newcommand{\AdfTasaNivTTp}{0{,}619}
\newcommand{\AdfTasaTend}{\ensuremath{-}1{,}00}
\newcommand{\AdfVCtm}{\ensuremath{-}2{,}87}
\newcommand{\AdfVCtt}{\ensuremath{-}3{,}42}
\newcommand{\AdfVivDif}{\ensuremath{-}5{,}30}
\newcommand{\AdfVivEst}{\ensuremath{-}3{,}65}
\newcommand{\AdfVivEstSBC}{\ensuremath{-}2{,}83}
\newcommand{\AdfVivEstSBCp}{14}
\newcommand{\AdfVivEstSBCpv}{0{,}055}
\newcommand{\AdfVivEstVC}{\ensuremath{-}2{,}87}
\newcommand{\AdfVivEstp}{0{,}005}
\newcommand{\AdfVivNivTM}{\ensuremath{-}1{,}86}
\newcommand{\AdfVivNivTMp}{0{,}352}
\newcommand{\AdfVivNivTT}{\ensuremath{-}1{,}79}
\newcommand{\AdfVivNivTTp}{0{,}707}
\newcommand{\AdfVivTend}{0{,}07}
\newcommand{\PacTasaDifDos}{\ensuremath{-}0{,}15}
\newcommand{\RhoDifDoce}{0{,}53}
\newcommand{\RhoDifTreintayseis}{0{,}48}
\newcommand{\RhoDifVeinticuatro}{0{,}48}
\newcommand{\RhoFdosDoce}{0{,}19}
\newcommand{\RhoFdosUno}{0{,}77}
\newcommand{\RhoFdosVeinticuatro}{0{,}18}
\newcommand{\RhoFtresDoce}{\ensuremath{-}0{,}45}
\newcommand{\RhoFtresOnce}{0{,}17}
\newcommand{\RhoFtresTrece}{0{,}20}
\newcommand{\RhoFtresUno}{\ensuremath{-}0{,}39}
\newcommand{\RhoFtresVeinticuatro}{\ensuremath{-}0{,}04}
\newcommand{\RhoNivUno}{0{,}93}
\newcommand{\RhoNivVeinticuatro}{0{,}66}
\newcommand{\RhoTasaDifDos}{\ensuremath{-}0{,}01}
\newcommand{\RhoTasaDifUno}{0{,}35}
\newcommand{\RhoTasaNivUno}{0{,}99}
\newcommand{\RhoTasaNivVeinticuatro}{0{,}73}

\newcommand{\AOV}{\ensuremath{-}31{,}8}
\newcommand{\AOVee}{6{,}8}
\newcommand{\AOVpct}{\ensuremath{-}27{,}3}
\newcommand{\AtipicosV}{ninguno}
\newcommand{\FiltroDSBCdos}{9{,}2}
\newcommand{\FiltroDSBCdost}{15{,}3}
\newcommand{\FiltroDSBCuno}{49{,}1}
\newcommand{\FiltroDosRaizAR}{1{,}012}
\newcommand{\FiltroUnoModelo}{ARMA(0,1)(1,1)\ensuremath{_{12}}}
\newcommand{\JBV}{4{,}0}
\newcommand{\JBVp}{0{,}133}
\newcommand{\JBVsinAO}{40}
\newcommand{\LBVdoce}{0{,}41}
\newcommand{\LBVtreintayseis}{0{,}53}
\newcommand{\LBVveinticuatro}{0{,}47}
\newcommand{\LBcuadV}{\ensuremath{<0{,}001}}
\newcommand{\NV}{487}
\newcommand{\RaizMAV}{2{,}23}
\newcommand{\RaizSMAV}{1{,}015}
\newcommand{\RejillaAdm}{20}
\newcommand{\RejillaN}{48}
\newcommand{\ResidAbrSinAO}{\ensuremath{-}5{,}1}
\newcommand{\RhoSAdoce}{\ensuremath{-}0{,}16}
\newcommand{\RhoSAveinticuatro}{\ensuremath{-}0{,}13}
\newcommand{\SAlbdoce}{0{,}62}
\newcommand{\SAlbtreintayseis}{0{,}88}
\newcommand{\SAlbveinticuatro}{0{,}65}
\newcommand{\SAmodelo}{ARMA(1,2)(1,1)\ensuremath{_{12}}}
\newcommand{\SAsigma}{6{,}78}
\newcommand{\SThetaSubDos}{0{,}81}
\newcommand{\SThetaSubUno}{0{,}80}
\newcommand{\SThetaV}{0{,}838}
\newcommand{\SThetaVee}{0{,}028}
\newcommand{\SigmaSubDos}{9{,}27}
\newcommand{\SigmaSubUno}{7{,}46}
\newcommand{\SigmaV}{8{,}25}
\newcommand{\SobreMaxT}{1{,}32}
\newcommand{\ThetaSubDos}{0{,}42}
\newcommand{\ThetaSubUno}{0{,}49}
\newcommand{\ThetaV}{0{,}449}
\newcommand{\ThetaVee}{0{,}039}

\newcommand{\AtipicosT}{2008-12 (\ensuremath{-}3{,}9); 2013-06 (3{,}9); 2022-04 (4{,}5); 2022-09 (5{,}2)}
\newcommand{\ConstT}{sin constante}
\newcommand{\CurtosisT}{6{,}0}
\newcommand{\JBT}{246}
\newcommand{\JBTp}{\ensuremath{<0{,}001}}
\newcommand{\LBTdoce}{0{,}93}
\newcommand{\LBTtreintayseis}{0{,}37}
\newcommand{\LBTveinticuatro}{0{,}32}
\newcommand{\LBcuadT}{\ensuremath{<0{,}001}}
\newcommand{\NT}{499}
\newcommand{\RhoTeoT}{0{,}37}
\newcommand{\SigmaT}{3{,}13}
\newcommand{\ThetaT}{\ensuremath{-}0{,}439}
\newcommand{\ThetaTee}{0{,}041}

\newcommand{\CoberturaAbanico}{100}
\newcommand{\MapeFtresDoce}{8{,}6}
\newcommand{\MapeFtresUno}{5{,}1}
\newcommand{\MejoraDoce}{31}
\newcommand{\MejoraSeis}{31}
\newcommand{\MejoraTasaDoce}{0}
\newcommand{\MejoraTasaUno}{10}
\newcommand{\MejoraTres}{34}
\newcommand{\MejoraUno}{19}
\newcommand{\NEval}{24}
\newcommand{\OrdenFdos}{ARMA(1,1)(0,1)\ensuremath{_{12}}}
\newcommand{\OrdenFtres}{ARMA(0,1)(0,1)\ensuremath{_{12}}}
\newcommand{\OrdenTasa}{ARMA(0,1)}
\newcommand{\RmseFtresDoce}{11{,}7}
\newcommand{\RmseIngDoce}{11{,}9}

\newcommand{\BandaCC}{0{,}089}
\newcommand{\CCsig}{$k=-24$ (\ensuremath{-}0{,}10), $k=-16$ (\ensuremath{-}0{,}11), $k=-6$ (\ensuremath{-}0{,}13), $k=-5$ (\ensuremath{-}0{,}11), $k=-3$ (\ensuremath{-}0{,}09), $k=1$ (0{,}15), $k=18$ (0{,}10)}
\newcommand{\DumDW}{0{,}092}
\newcommand{\DumR}{0{,}255}
\newcommand{\ElasDoce}{\ensuremath{-}0{,}54}
\newcommand{\ElasDoceee}{0{,}33}
\newcommand{\ElasTres}{\ensuremath{-}0{,}35}
\newcommand{\ElasTresee}{0{,}20}
\newcommand{\FilDW}{2{,}82}
\newcommand{\FilR}{0{,}019}
\newcommand{\HaughDoceNeg}{0{,}014}
\newcommand{\HaughDocePos}{0{,}147}
\newcommand{\HaughSeisP}{\ensuremath{<0{,}001}}
\newcommand{\NivDW}{0{,}149}
\newcommand{\NivPend}{0{,}31}
\newcommand{\NivR}{0{,}107}
\newcommand{\NivT}{7{,}7}
\newcommand{\PPDoce}{\ensuremath{-}7{,}6}
\newcommand{\PPDoceHi}{2}
\newcommand{\PPDoceLo}{\ensuremath{-}17}
\newcommand{\PPequiv}{14{,}0}
\newcommand{\RDdw}{2{,}84}
\newcommand{\RDrdos}{0{,}065}
\newcommand{\TasaUltima}{6{,}67}


\newcommand{\nota}[2]{\textit{(Nota~#1, Sec.~#2)}}
\newtcolorbox{clave}{colback=blue!4,colframe=blue!45!black,boxrule=0.6pt,arc=2pt,left=6pt,right=6pt,top=4pt,bottom=4pt,title={Idea clave},fonttitle=\bfseries\small}
\newtcolorbox{pregunta}[1]{colback=orange!5,colframe=orange!60!black,boxrule=0.5pt,arc=2pt,left=6pt,right=6pt,top=3pt,bottom=3pt,title={#1},fonttitle=\bfseries\small}
\newcommand{\Q}[2]{\begin{pregunta}{#1}#2\end{pregunta}}

\title{\textbf{Gu\'ia para defender el TP}\\[2pt]\large Inicios de viviendas y tasa hipotecaria: qu\'e hice, por qu\'e, c\'omo se lee y qu\'e me pueden preguntar}
\author{Econometr\'ia II --- C\'atedra Mirta L.\ Gonz\'alez}
\date{Octubre de 2026}

\begin{document}
\maketitle
\tableofcontents
\newpage

%======================================================================
\section{C\'omo usar esta gu\'ia}
Cada secci\'on sigue el mismo orden que el informe y tiene cuatro partes:
\begin{itemize}
\item \textbf{Qu\'e hice}: el paso en dos l\'ineas.
\item \textbf{Por qu\'e}: la teor\'ia que lo respalda, con el \'algebra que hace falta para poder justificarlo en el pizarr\'on, y la cita a la Nota.
\item \textbf{C\'omo se lee}: qu\'e mirar en la salida y qu\'e n\'umero decir.
\item \textbf{Preguntas probables} (recuadros naranja) con la respuesta corta, y una \textbf{idea clave} (recuadro azul) para cerrar.
\end{itemize}
Los n\'umeros salen del mismo c\'odigo que el informe, as\'i que coinciden con \'el. Al final hay una hoja con los n\'umeros que conviene saber de memoria (Secci\'on~\ref{sec:numeros}).

%======================================================================
\section{La historia en dos minutos (guion oral)}
\begin{enumerate}[leftmargin=1.4em]
\item \textbf{Pregunta.} Quiero medir cu\'anto y con qu\'e demora caen los inicios de construcci\'on de viviendas en EE.UU.\ cuando sube la tasa hipotecaria. La tasa es la causa: encarece la cuota y frena la decisi\'on de empezar una obra. No es parte de los inicios, y se forma en el mercado de bonos.
\item \textbf{Datos.} Mensuales, 1985--2026, 500 observaciones. Arranco en 1985 porque en la \'epoca de Volcker la tasa era 2,6 veces m\'as vol\'atil y eso rompe la varianza constante.
\item \textbf{Transformaci\'on.} Trabajo en logaritmos: la relaci\'on rango--mediana da $\lambda$ cerca de cero en las dos series.
\item \textbf{Inicios.} Tienen estacionalidad fuerte y cambiante y dos ra\'ices unitarias, una regular y una estacional. El correlograma lo muestra y el ADF lo confirma para la regular. Con $(1-B)(1-B^{12})$ el correlograma tiene la firma del modelo de l\'ineas a\'ereas, y eso es lo que estimo: $\hat\theta=\ThetaV$, $\hat\Theta=\SThetaV$, con un at\'ipico en abril de 2020. Los residuos son ruido blanco.
\item \textbf{Tasa.} Una ra\'iz unitaria, sin estacionalidad. En diferencias es un MA(1).
\item \textbf{Pron\'ostico.} Fuera de muestra el modelo de los inicios le gana al paseo aleatorio estacional entre un 19 y un 34\%.
\item \textbf{Relaci\'on.} En niveles la regresi\'on es espuria: da que m\'as tasa son m\'as inicios. Filtrando igual las dos series, la elasticidad a un a\~no es $\ElasDoce$: unos $\PPDoce\%$ de inicios por punto de tasa, pero imprecisa. Adem\'as hay retroalimentaci\'on: los inicios tambi\'en anticipan a la tasa. Eso es lo que hay que resolver en la segunda parte con la funci\'on de transferencia.
\end{enumerate}

%======================================================================
\section{Las reglas de la c\'atedra y d\'onde se cumplen}
\begin{center}\small
\begin{tabularx}{\linewidth}{>{\raggedright}p{6.2cm}X}
\toprule
Regla (clase del 14/10) & D\'onde est\'a en el informe \\
\midrule
1. Objetivo antes que t\'ecnica & Secci\'on 1: la pregunta y el mecanismo \\
2. La segunda serie es una causa, no una parte & Secci\'on 1 (por qu\'e la tasa es causa) y 8 (retroalimentaci\'on) \\
3. M\'as de 30 observaciones, mejor 50 o m\'as & 500 meses \\
4. Primero cada serie, despu\'es la relaci\'on & Secciones 2--6 antes de la 8 \\
5. Graficar con contexto & Figura 1, hitos hist\'oricos, at\'ipicos con fecha y explicaci\'on \\
6. El correlograma manda; FAC gruesa o AR(3)--AR(4) $\Rightarrow$ dudar & FAC de $(1-B^{12})x$, ra\'iz AR de 1,012 en F2 \\
7. Dickey--Fuller no es receta & ADF con 12, 24 y SBC rezagos; ADF fr\'agil sobre $(1-B^{12})x$ \\
8. Testear estacionalidad; probar serie con y sin desestacionalizar & $F$ de dummies y de cambio de patr\'on; modelo de la serie desestacionalizada \\
9. Cambios estructurales por subper\'iodos & Muestra desde 1985; modelos 1986--2007 vs.\ 2008--2026 \\
10. Residuos ruido blanco no alcanza; parsimonia & Rejilla con SBC, sobreajuste, cuadrados, Jarque--Bera \\
11. Lo que se eval\'ua es el criterio & Secci\'on 9: limitaciones dichas expl\'icitamente \\
12. Software que muestre todo & R, con salidas en formato EViews (correlograma, $Q$, ra\'ices, residuos) \\
\bottomrule
\end{tabularx}
\end{center}

%======================================================================
\section{Datos y muestra}
\textbf{Qu\'e hice.} Inicios sin desestacionalizar (HOUSTNSA, Census) y tasa fija a 30 a\~nos (MORTGAGE30US, Freddie Mac). La tasa es semanal: la paso a mensual con el promedio de las semanas del mes, que es lo que hace FRED. Muestra principal: 1985-01 a 2026-08.

\textbf{Por qu\'e la serie sin desestacionalizar.} La serie desestacionalizada ya pas\'o por X-13, un filtro externo que no controlo. Si la modelo, estoy modelando la serie m\'as el filtro. Con la serie original la estacionalidad se trata dentro del modelo con $(1-B^{12})$ y un SMA, que es lo que ense\~na la materia. Igual la prob\'e (regla 8; Secci\'on~\ref{sec:sa}).

\textbf{Por qu\'e desde 1985.} El desv\'io de la variaci\'on mensual de la tasa es $\VolVolcker$\,pp en 1979--84 contra $\VolNormal$\,pp en 1985--2007. Un modelo ARMA supone $\sigma^2_a$ constante. Mezclar los dos reg\'imenes har\'ia que el modelo promedie dos varianzas distintas, y los diagn\'osticos fallar\'ian por eso y no por la din\'amica.

\Q{¿Por qu\'e no usar toda la serie desde 1971?}{Porque el r\'egimen de Volcker es un cambio estructural en la varianza de la tasa (2,6 veces), y la regla 9 pide trabajar por subper\'iodos. Ganar 14 a\~nos de datos no compensa violar el supuesto de varianza constante. Con 500 observaciones sobra muestra.}

\Q{¿Por qu\'e promedio y no el dato de fin de mes?}{Los inicios son un flujo de todo el mes: las obras arrancan a lo largo del mes con la tasa vigente en cada semana. El promedio representa mejor la tasa que enfrent\'o el mes.}

%======================================================================
\section{Transformaci\'on: por qu\'e logaritmos \nota{6}{5.4}}
\textbf{Qu\'e hice.} Para cada a\~no calcul\'e la mediana $M_i$ y el rango $R_i$ y regres\'e $\ln R_i$ sobre $\ln M_i$. Como contraste, el perfil de verosimilitud de Box--Cox.

\textbf{Por qu\'e (\'algebra).} Si el desv\'io crece con el nivel como $\sigma\propto\mu^{1-\lambda}$, la transformaci\'on $T(Y)=Y^\lambda$ lo estabiliza. Por el m\'etodo delta, $\operatorname{Var}[T(Y)]\approx[T'(\mu)]^2\sigma^2=\lambda^2\mu^{2(\lambda-1)}\sigma^2$, que es constante cuando $\sigma\propto\mu^{1-\lambda}$. El rango es proporcional al desv\'io, as\'i que
\[
\ln R_i = a + (1-\lambda)\ln M_i \quad\Longrightarrow\quad \hat\lambda = 1-\hat b .
\]
Inicios: $\hat b=0{,}78\Rightarrow\hat\lambda=\LamInicios$ (IC 95\%: $[\LamInicioslo;\LamIniciosup]$). Tasa: $\hat\lambda=\LamTasa$ ($[\LamTasalo;\LamTasaup]$). Los dos intervalos contienen el 0 (logaritmo) y excluyen el 1 (no transformar).

\textbf{El perfil de Box--Cox y el jacobiano.} Para comparar verosimilitudes con distintos $\lambda$ hay que sumar el jacobiano: si $z=T(y)$, la densidad de $y$ es $f_z(T(y))\cdot|T'(y)|$, y con $T(y)=100(y^\lambda-1)/\lambda$ queda $\ell_y(\lambda)=\ell_z(\lambda)+(\lambda-1)\sum\ln y_t+\text{cte}$. Sin ese t\'ermino se comparar\'ian verosimilitudes de variables distintas. El perfil tiene el m\'aximo en $\lambda=1/2$, pero tambi\'en descarta $\lambda=1$.

\textbf{Por qu\'e igual me quedo con el logaritmo.} (i)~Es lo que da el m\'etodo del apunte. (ii)~Las diferencias de $100\ln$ son variaciones porcentuales y la relaci\'on se lee como elasticidad.

\Q{Si la verosimilitud prefiere ra\'iz cuadrada, ¿por qu\'e logaritmo?}{Porque el criterio del apunte (rango--mediana) es compatible con el logaritmo y la interpretaci\'on es mucho m\'as clara (porcentajes y elasticidades). El perfil de verosimilitud es sensible a las colas, como la ca\'ida de abril de 2020.}

\Q{¿Por qu\'e multiplic\'as el logaritmo por 100?}{Solo por escala: $100\,\Delta\ln x_t\approx$ variaci\'on porcentual. As\'i $\hat\sigma_a=\SigmaV$ se lee directamente como ``el error de pron\'ostico a un mes es de un 8\%''. No cambia ning\'un test.}

\begin{clave}La transformaci\'on se decide \emph{antes} que la diferenciaci\'on: diferenciar no corrige la heterocedasticidad. Primero se estabiliza la varianza, despu\'es se discuten las ra\'ices unitarias.\end{clave}

%======================================================================
\section{Estacionalidad}
\textbf{Qu\'e hice.} Regres\'e $\Delta x_t$ sobre 11 dummies mensuales centradas: $F$ de estacionalidad fija. Despu\'es agregu\'e la interacci\'on mes $\times$ subper\'iodo (1985--2007 vs.\ 2008--2026): $F$ de cambio de patr\'on.

\textbf{Por qu\'e dos tests.} Hay dos tipos de estacionalidad, y cada uno se trata distinto:
\begin{itemize}
\item \textbf{Determinista}: el patr\'on es fijo, $x_t=\sum_j\delta_jD_{jt}+\dots$ Se trata con dummies.
\item \textbf{Estoc\'astica}: hay ra\'ices unitarias estacionales y el patr\'on \emph{migra} con el tiempo. Se trata con $(1-B^{12})$.
\end{itemize}
El primer $F$ dice si hay estacionalidad ($F=\EstFv$, $R^2=\EstRv$). El segundo dice de qu\'e tipo es: si el patr\'on cambia entre subper\'iodos ($F=\EstFcv$, $p\ \EstPcv$), no es fijo, as\'i que es estoc\'astica. La tasa no tiene estacionalidad ($p=\EstPt$).

\textbf{\'Algebra de $(1-B^{12})$.} $1-B^{12}=(1-B)(1+B+B^2+\dots+B^{11})$, y sus 12 ra\'ices son las ra\'ices duod\'ecimas de la unidad, $z_k=e^{i2\pi k/12}$. Una es $z=1$ (la ra\'iz regular) y las otras 11 son estacionales \nota{6}{2.3}. Por eso $(1-B^{12})$ ya contiene un $(1-B)$, y aplicar $(1-B)(1-B^{12})$ significa que hay \emph{dos} ra\'ices en $z=1$.

\Q{¿Por qu\'e dummies centradas?}{Cada dummy es ``mes $j$ menos diciembre'', as\'i suman cero en el a\~no y la constante sigue siendo el promedio anual. Adem\'as, con dummies centradas la distribuci\'on del estad\'istico $\tau$ del ADF no cambia (Dickey, Bell y Miller), y eso me deja usarlas dentro del ADF.}

\Q{¿Por qu\'e no hiciste HEGY?}{Es el test espec\'ifico de ra\'ices estacionales, pero est\'a fuera del programa (la Nota 6 lo menciona sin desarrollarlo). Lo reemplac\'e con evidencia indirecta: el $F$ de cambio de patr\'on, la FAC de $\Delta x$ en 12, 24 y 36 y la comparaci\'on de filtros por SBC.}

%======================================================================
\section{Estacionariedad: correlogramas y Dickey--Fuller}

\subsection{C\'omo leer los correlogramas \nota{6}{9}}
\textbf{La banda.} Bajo ruido blanco $\hat\rho_k\sim N(0,1/T)$ (Bartlett) y $\hat\phi_{kk}\sim N(0,1/T)$ (Quenouille), de ah\'i la banda $\pm1{,}96/\sqrt T\approx\pm0{,}088$ con $T=500$. Con 36 rezagos es esperable que 1 o 2 caigan afuera por azar; por eso se mira tambi\'en el $Q$ conjunto.

\textbf{El $Q$ de Ljung--Box.} $Q_K=T(T+2)\sum_{k=1}^K\hat\rho_k^2/(T-k)\sim\chi^2_K$ bajo $H_0$: $\rho_1=\dots=\rho_K=0$. Sobre residuos de un modelo con $k$ par\'ametros ARMA los grados de libertad son $K-k$. En la salida tipo EViews, ``Prob'' es el $p$-valor de ese $Q$ acumulado hasta el rezago de la fila.

\textbf{Por qu\'e una ra\'iz unitaria da una FAC que no cae.} En una caminata aleatoria $\rho(Y_t,Y_{t-k})=\sqrt{(t-k)/t}\to1$ para todo $k$ fijo \nota{6}{3.1}. La FAC muestral parte cerca de 1 y decae casi linealmente, que es lo que se ve en el nivel de las dos series ($\RhoNivUno$ y $\RhoTasaNivUno$ en $k=1$).

\textbf{Los cuatro filtros de los inicios, uno por uno:}
\begin{enumerate}[leftmargin=1.4em]
\item \textbf{Nivel:} decae lento y tiene picos en 12 y 24. Hay ra\'iz regular y estacional.
\item \textbf{$(1-B)x$:} se va la persistencia de corto plazo, pero $\rho_{12}=\RhoDifDoce$, $\rho_{24}=\RhoDifVeinticuatro$, $\rho_{36}=\RhoDifTreintayseis$ no caen. Es la versi\'on estacional del mismo patr\'on: una ra\'iz unitaria en las frecuencias estacionales.
\item \textbf{$(1-B^{12})x$:} los picos estacionales desaparecen, pero la FAC es ``gruesa'' y cae despacio desde $\RhoFdosUno$. Seg\'un la regla 6, falta la diferencia regular.
\item \textbf{$(1-B)(1-B^{12})x$:} solo quedan $\rho_1=\RhoFtresUno$, $\rho_{12}=\RhoFtresDoce$ y los ``sat\'elites'' $\rho_{11}$ y $\rho_{13}$ positivos. Es la firma del modelo de l\'ineas a\'ereas (Secci\'on~\ref{sec:airline}).
\end{enumerate}

\Q{¿C\'omo s\'e que no sobrediferenci\'e?}{Si una serie estacionaria se diferencia de m\'as, aparece una ra\'iz unitaria en el MA y $\rho_1\to-0{,}5$. Por ejemplo, si $Y_t=\varepsilon_t$ y diferencio, $\Delta Y_t=\varepsilon_t-\varepsilon_{t-1}$ tiene $\rho_1=-1/2$ \nota{6}{2.6}. Ac\'a $\rho_1=\RhoFtresUno$ y $\hat\theta=\ThetaV$, lejos de 1. Para la parte estacional, ver la pregunta sobre $\Theta$ en la Secci\'on~\ref{sec:airline}.}

\subsection{Dickey--Fuller aumentado \nota{6}{4}}
\textbf{La regresi\'on.}
\[
\Delta x_t=\alpha_0+\alpha_1t+\gamma\,x_{t-1}+\sum_{j=1}^{p}c_j\,\Delta x_{t-j}+\text{dummies}+e_t,\qquad H_0:\gamma=0\ \text{(ra\'iz unitaria)},\quad H_1:\gamma<0 .
\]
Sale de reparametrizar un AR($p+1$): $\gamma=\sum\phi_i-1$, as\'i que $\gamma=0$ equivale a $\Phi(1)=0$, una ra\'iz en $z=1$ \nota{6}{4.4}. Los rezagos $\Delta x_{t-j}$ est\'an para que el residuo sea ruido blanco.

\textbf{Por qu\'e no se usa la $t$ de Student.} Bajo $H_0$, $x_{t-1}$ es $I(1)$ y su varianza crece con $t$. La $t$ de $\hat\gamma$ converge a un funcional del movimiento browniano, no a una normal. Por eso los valores cr\'iticos son los de Dickey--Fuller (m\'as negativos) y por eso en la salida la columna ``Prob.'' de la fila del nivel rezagado \emph{no sirve}: el $p$-valor v\'alido es el de arriba, ``Prob.*'' de MacKinnon.

\textbf{Las tres variantes y el orden secuencial.} $\tau$ (sin nada), $\tau_\mu$ (constante), $\tau_\tau$ (constante y tendencia), cada una con su tabla. Se arranca por la m\'as general ($\tau_\tau$). Si la tendencia no es significativa se baja a $\tau_\mu$, porque incluir una tendencia que no est\'a le quita potencia al test \nota{6}{4.3}.

\textbf{El esquema descendente (Dickey--Pantula) \nota{6}{4.5}.} Si hubiera dos ra\'ices, el ADF sobre el nivel no es v\'alido, porque supone que $\Delta x$ es estacionaria. Por eso:
\begin{enumerate}[leftmargin=1.4em]
\item \textbf{Paso 1}: ADF sobre $\Delta x$ (¿hay una segunda ra\'iz?). Inicios: $\AdfVivDif$; tasa: $\AdfTasaDif$, contra $\AdfVCtm$. Se rechaza: hay a lo sumo una ra\'iz.
\item \textbf{Paso 2}: ADF sobre el nivel. Inicios: $\tau_\tau=\AdfVivNivTT$ ($p=\AdfVivNivTTp$), tendencia con $t=\AdfVivTend$, as\'i que bajo a $\tau_\mu=\AdfVivNivTM$ ($p=\AdfVivNivTMp$). Tasa: $\tau_\mu=\AdfTasaNivTM$ ($p=\AdfTasaNivTMp$). No se rechaza: hay una ra\'iz. Las dos series son $I(1)$.
\end{enumerate}

\textbf{¿Cu\'antos rezagos?} Said--Dickey: un ARIMA($p,1,q$) se aproxima por un AR largo si $p^*\le n^{1/3}\approx8$ \nota{6}{4.6}. Con datos mensuales estacionales tom\'e 12 (un a\~no). Como sensibilidad, 24 rezagos y el $p$ que minimiza el SBC: la conclusi\'on no cambia.

\Q{Si el ADF sobre $\Delta x$ rechaza, ¿no es estacionaria $\Delta x$? ¿Para qu\'e diferenci\'as estacionalmente?}{Porque el ADF solo mira la ra\'iz $z=1$: regresa $\Delta^2x_t$ sobre $\Delta x_{t-1}$, y eso no ve las ra\'ices estacionales $e^{i2\pi k/12}$ \nota{6}{2.3}. Que rechace solo dice que no hay \emph{otra} ra\'iz en 1. Las ra\'ices estacionales las muestran la FAC de $\Delta x$ (picos en 12, 24 y 36 que no caen) y el $F$ de cambio de patr\'on.}

\Q{El ADF sobre $(1-B^{12})x$ rechaza con 12 rezagos. ¿Por qu\'e igual aplic\'as $(1-B)$?}{Porque el resultado es fr\'agil: con el $p$ que elige el SBC ($\AdfVivEstSBCp$ rezagos) da $p=\AdfVivEstSBCpv$. Y tres cosas apuntan para el otro lado: (i)~la FAC de $(1-B^{12})x$ es gruesa (regla 6); (ii)~el mejor modelo sobre $(1-B^{12})x$ necesita un AR(1) con ra\'iz de m\'odulo $\FiltroDosRaizAR$, o sea $\hat\phi\approx0{,}99$, que es la ra\'iz unitaria disfrazada (regla 6: ``tener que poner un AR deber\'ia hacerte dudar''); (iii)~la comparaci\'on por SBC da ganador al doble filtro. Es justo la regla 7: el DF no es una receta.}

\begin{clave}El ADF responde una sola pregunta: ¿hay una ra\'iz en $z=1$? Las ra\'ices estacionales y la elecci\'on entre filtros las decide el correlograma, con el ADF como un elemento m\'as.\end{clave}

\subsection{Comparar filtros con el SBC}
\textbf{Qu\'e hice.} Para cuatro filtros (F1: $(1-B)$ con dummies; F2: $(1-B^{12})$; F2t: $(1-B^{12})$ con tendencia; F3: $(1-B)(1-B^{12})$) busqu\'e el mejor ARMA admisible y compar\'e los SBC, todos sobre las mismas fechas (desde 1986-02).

\textbf{Por qu\'e se pueden comparar.} Los cuatro son modelos para la \emph{misma} variable $x_t$. Si $W_t=F(B)x_t$ con $F(0)=1$, entonces $W_t=x_t+(\text{combinaci\'on de valores pasados de }x)$. Dados los valores iniciales, el paso de $(x_t)$ a $(W_t)$ es una transformaci\'on triangular con unos en la diagonal: su jacobiano es 1. La verosimilitud de $W$ es la verosimilitud (condicional) de $x$, y los SBC son comparables. Es distinto de comparar $\ln y$ con $y$, donde s\'i hay que sumar el jacobiano. Resultado: F3 gana por $\FiltroDSBCdos$ puntos frente a F2 y por $\FiltroDSBCuno$ frente a las dummies.

\Q{¿Qu\'e significa una diferencia de 9 puntos de SBC?}{Por la convenci\'on usual, una diferencia menor a 2 es evidencia d\'ebil y conviene combinar o mirar otros criterios \nota{6}{11.5}. Nueve puntos es una diferencia clara a favor del doble filtro.}

%======================================================================
\section{Inicios: el modelo de l\'ineas a\'ereas}\label{sec:airline}

\subsection{El modelo y su FAC te\'orica (Nota 4)}
\[
W_t\equiv(1-B)(1-B^{12})x_t=(1-\theta B)(1-\Theta B^{12})a_t=a_t-\theta a_{t-1}-\Theta a_{t-12}+\theta\Theta\,a_{t-13}.
\]
Es un MA(13) con solo dos par\'ametros libres. Su FAC se calcula como la de cualquier MA, multiplicando los coeficientes que coinciden:
\begin{align*}
\gamma_0&=\sigma^2(1+\theta^2+\Theta^2+\theta^2\Theta^2)=\sigma^2(1+\theta^2)(1+\Theta^2),\\
\gamma_1&=\sigma^2[(-\theta)(1)+(-\Theta)(\theta\Theta)]=-\theta\sigma^2(1+\Theta^2)\ \Rightarrow\ \rho_1=\frac{-\theta}{1+\theta^2},\\
\gamma_{12}&=\sigma^2[(-\Theta)(1)+(\theta\Theta)(-\theta)]=-\Theta\sigma^2(1+\theta^2)\ \Rightarrow\ \rho_{12}=\frac{-\Theta}{1+\Theta^2},\\
\gamma_{11}&=\gamma_{13}=\sigma^2\theta\Theta\ \Rightarrow\ \rho_{11}=\rho_{13}=\frac{\theta\Theta}{(1+\theta^2)(1+\Theta^2)},
\end{align*}
y $\rho_k=0$ para el resto. Con $\hat\theta=\ThetaV$ y $\hat\Theta=\SThetaV$: $\rho_1=-0{,}37$ (emp\'irica $\RhoFtresUno$), $\rho_{12}=-0{,}49$ (emp\'irica $\RhoFtresDoce$), $\rho_{11}=\rho_{13}=0{,}18$ (emp\'iricas $\RhoFtresOnce$ y $\RhoFtresTrece$). Los ``sat\'elites'' positivos en 11 y 13 son la huella del producto $\theta\Theta$: por eso el modelo es \emph{multiplicativo}.

\textbf{Interpretaci\'on.} $\theta$: parte de cada sorpresa mensual se corrige al mes siguiente. $\Theta=0{,}84$: el patr\'on estacional se mueve, pero despacio; cada a\~no se ajusta un poco seg\'un la sorpresa del mismo mes del a\~no anterior.

\textbf{Por qu\'e se llama de l\'ineas a\'ereas.} Box y Jenkins lo usaron con los pasajeros a\'ereos internacionales. Es el modelo est\'andar para series mensuales con tendencia estoc\'astica y estacionalidad estoc\'astica.

\subsection{Identificaci\'on: la rejilla \nota{6}{7--8}}
Estim\'e los $\RejillaN$ ARMA$(p\le3,q\le2)(P\le1,Q\le1)_{12}$ sobre $W_t$. Un modelo es \textbf{admisible} si es estacionario (ra\'ices AR fuera del c\'irculo), invertible (ra\'ices MA fuera) y sus residuos pasan Ljung--Box en 12, 24 y 36. Entre los admisibles ($\RejillaAdm$) elijo el de menor SBC. El ganador coincide con lo que dec\'ia el correlograma. El segundo, ARMA(2,0)(0,1), tiene 3 par\'ametros y 3,4 puntos m\'as de SBC: un AR(2) que imita al MA(1), con un par\'ametro de m\'as (parsimonia, \nota{6}{7}).

\Q{¿Por qu\'e SBC y no AIC?}{El SBC penaliza m\'as ($\ln T\approx6{,}2$ por par\'ametro contra 2) y es consistente: si el modelo verdadero es un ARMA finito, lo encuentra con probabilidad que tiende a 1. El AIC tiende a sobreparametrizar \nota{6}{8.3}. Ac\'a el AIC elige un modelo casi empatado (ARMA(2,0)), as\'i que da lo mismo.}

\subsection{La salida: c\'omo se lee l\'inea por l\'inea}
\begin{itemize}
\item \textbf{Dependent Variable: D(LHOUST,1,12)}: es $W_t=(1-B)(1-B^{12})\,100\ln x_t$. Las 13 primeras observaciones se pierden al diferenciar: la muestra arranca en 1986M02 y tiene 487 datos.
\item \textbf{MA(1) $=-0{,}4487$, SMA(12) $=-0{,}8384$}: EViews y R escriben $(1+\theta B)$, las Notas escriben $(1-\theta B)$. Por eso en las Notas $\hat\theta=+\ThetaV$ y $\hat\Theta=+\SThetaV$. Errores est\'andar chicos ($t$ de $-11{,}6$ y $-29{,}7$).
\item \textbf{AO2020M04 $=\AOV$}: el efecto del at\'ipico, en unidades de $100\ln$. En porcentaje: $e^{-0{,}318}-1=\AOVpct\%$.
\item \textbf{R-squared 0,53}: la parte de la varianza de $W_t$ que explica el modelo. Con series diferenciadas el $R^2$ es bajo por naturaleza y no es un criterio.
\item \textbf{S.E.\ of regression $\approx8{,}3$}: desv\'io del error a un paso, en \%.
\item \textbf{Akaike / Schwarz / Hannan--Quinn}: por observaci\'on, $\text{SC}=-2\ell/n+k\ln n/n$. Solo sirven para comparar modelos sobre la misma muestra.
\item \textbf{Durbin--Watson 2,05}: cerca de 2, sin autocorrelaci\'on de orden 1 en los residuos. En un ARMA es informativo, pero el test formal es el Ljung--Box.
\item \textbf{Jarque--Bera 4,0 ($p=0{,}13$)}: no se rechaza normalidad.
\item \textbf{Inverted MA Roots}: las ra\'ices de $\Theta(z)=0$ invertidas, $1/z$. El modelo es invertible si todas tienen m\'odulo $<1$. Hay 0,45 (es $\hat\theta$) y doce ra\'ices de m\'odulo $0{,}985$, que vienen de $1-\Theta z^{12}=0$, cuyas ra\'ices cumplen $|z|=\Theta^{-1/12}$. Sus inversas tienen m\'odulo $\Theta^{1/12}=0{,}838^{1/12}=0{,}985$, distribuidas en los 12 \'angulos del c\'irculo. Por eso aparecen $0{,}99$, $-0{,}99$, $0{,}85\pm0{,}49i$, etc.
\end{itemize}

\Q{$\Theta=0{,}84$ est\'a cerca de 1. ¿No es una se\~nal de sobrediferenciaci\'on estacional?}{Si $\Theta$ fuera 1, el factor $(1-B^{12})$ de los dos lados se cancelar\'ia: la estacionalidad ser\'ia determinista y bastar\'ian dummies. Pero (i)~el IC 95\% de $\hat\Theta$ es $\SThetaV\pm1{,}96\times\SThetaVee=[0{,}78;\ 0{,}89]$, que excluye el 1; (ii)~el modelo con dummies (F1) pierde por $\FiltroDSBCuno$ puntos de SBC; (iii)~el patr\'on estacional cambia entre subper\'iodos. La estacionalidad es estoc\'astica, aunque de movimiento lento.}

\Q{¿Por qu\'e un at\'ipico aditivo y no una dummy cualquiera?}{Un AO es un shock puntual que afecta solo ese mes: $x_t=\omega I_t^{2020.04}+N_t$, con $N_t$ el ARIMA. Como hay que filtrar los dos lados, el regresor entra como $(1-B)(1-B^{12})I_t$, que vale 1, $-1$, $-1$ y $1$ en 2020-04, 2020-05, 2021-04 y 2021-05. Hay contexto (cierre por COVID) y evidencia: sin el AO, ese mes es un residuo de $\ResidAbrSinAO$ desv\'ios y Jarque--Bera pasa de 4 a $\JBVsinAO$. Modelarlo como intervenci\'on con din\'amica es de la segunda parte.}

\subsection{Diagn\'ostico \nota{6}{10}}
\begin{itemize}
\item \textbf{Ljung--Box}: $p=\LBVdoce$, $\LBVveinticuatro$, $\LBVtreintayseis$ (12, 24 y 36 rezagos, $K-2$ gl). No queda autocorrelaci\'on.
\item \textbf{Jarque--Bera}: $\JBV$ ($p=\JBVp$): normalidad aceptable.
\item \textbf{Sobreajuste} \nota{6}{10.4}: agrego AR(1), MA(2), SAR(12) o SMA(24) de a uno. Ninguno es significativo ($|t|\le\SobreMaxT$) y todos empeoran el SBC. El modelo no est\'a subespecificado.
\item \textbf{Ljung--Box sobre los cuadrados} ($p\ \LBcuadV$): hay heterocedasticidad condicional, con m\'as volatilidad despu\'es de 2008. Es una limitaci\'on, y es exactamente lo que la regla 10 advierte: ruido blanco no alcanza.
\end{itemize}

\Q{¿Por qu\'e los grados de libertad son $K-2$?}{Porque los residuos salen de un modelo con 2 par\'ametros ARMA estimados ($\theta$ y $\Theta$), y estimarlos ``gasta'' 2 grados de libertad del $Q$. El coeficiente del AO no cuenta: es un regresor, no un par\'ametro de la din\'amica.}

\subsection{Estabilidad y serie desestacionalizada}\label{sec:sa}
Por subper\'iodos, $\hat\theta$ pasa de $\ThetaSubUno$ a $\ThetaSubDos$ y $\hat\Theta$ de $\SThetaSubUno$ a $\SThetaSubDos$: la \emph{din\'amica} es estable. Lo que cambia es $\hat\sigma$ ($\SigmaSubUno\to\SigmaSubDos$): los shocks son m\'as grandes despu\'es de 2008.

\textbf{Serie desestacionalizada (regla 8).} Su primera diferencia todav\'ia tiene $\rho_{12}=\RhoSAdoce$ y $\rho_{24}=\RhoSAveinticuatro$, y el mejor modelo admisible necesita t\'erminos estacionales ($\SAmodelo$). Es un efecto conocido: la desestacionalizaci\'on deja autocorrelaci\'on negativa en los rezagos estacionales. Desestacionalizar no me ahorra modelar la estacionalidad, y por eso uso la serie original.

\begin{clave}Dos par\'ametros explican toda la din\'amica de los inicios: $\theta$ (correcci\'on mes a mes) y $\Theta$ (movimiento lento del patr\'on estacional). La FAC te\'orica que generan es la que se ve en los datos.\end{clave}

%======================================================================
\section{Tasa hipotecaria: ARIMA(0,1,1)}
\textbf{Qu\'e hice.} $y_t=100\ln$ tasa, $I(1)$ sin estacionalidad. Sobre $(1-B)y_t$, rejilla ARMA($p\le3,q\le2$) con y sin constante. Gana:
\[
(1-B)y_t=(1-\theta B)a_t,\qquad\hat\theta=\ThetaT,\ \text{o sea}\ (1-B)y_t=a_t+0{,}44\,a_{t-1}.
\]

\textbf{C\'omo se lee.} $\theta$ es \emph{negativo} en la convenci\'on de las Notas, as\'i que $\rho_1=-\theta/(1+\theta^2)=+\RhoTeoT$: los cambios de tasa tienen inercia de un mes (si sube este mes, tiende a seguir subiendo el que viene). La FAC emp\'irica tiene un solo pico, $\rho_1=\RhoTasaDifUno$, y despu\'es se corta, como un MA(1) (Nota 4). La FACP alterna y decae, como corresponde a un MA.

\textbf{Sin constante.} Con constante habr\'ia una deriva: una tendencia determinista en el logaritmo de la tasa. La tasa pas\'o de 13\% a 6,7\% en 40 a\~nos, una ca\'ida de $0{,}13\%$ por mes en $100\ln$, chica frente a $\hat\sigma=\SigmaT$. El SBC la descarta (5,8 puntos peor).

\textbf{Diagn\'ostico.} Ljung--Box no rechaza ($p=\LBTdoce$, $\LBTveinticuatro$, $\LBTtreintayseis$), el sobreajuste no agrega nada y $\hat\theta$ es igual en los dos subper\'iodos. Pero los residuos tienen colas pesadas (curtosis $\CurtosisT$, Jarque--Bera $\JBT$) y volatilidad cambiante. Los residuos extremos tienen fecha y explicaci\'on:
\begin{itemize}
\item 2008-12: la Fed anuncia compras de hipotecas (MBS), y la tasa cae.
\item 2013-06: \emph{taper tantrum}, la Fed anuncia que va a reducir las compras y la tasa salta.
\item 2022-04 y 2022-09: el ciclo de suba de tasas m\'as r\'apido en 40 a\~nos.
\end{itemize}

\Q{Si los residuos no son normales, ¿sirve el modelo?}{Para la media condicional (el pron\'ostico puntual) s\'i: m\'axima verosimilitud gaussiana sigue dando estimadores consistentes aunque los errores no sean normales (cuasi-m\'axima verosimilitud). Lo que se ve afectado son los intervalos de pron\'ostico, que quedan angostos en las colas. La soluci\'on (ARCH/GARCH o intervenciones en esas fechas) est\'a fuera de la primera parte, y lo digo como limitaci\'on.}

%======================================================================
\section{Pron\'ostico fuera de muestra \nota{5}{8--9}}
\textbf{Dise\~no.} Estimo con 1985--2017 y eval\'uo en 2018--2019 (antes del COVID). El origen es m\'ovil: en cada uno de los $\NEval$ meses reestimo los par\'ametros y pronostico a 1, 3, 6 y 12 meses. El \emph{orden} del modelo lo eleg\'i solo con la muestra de estimaci\'on. Si lo eligiera mirando 2018--2019, la evaluaci\'on estar\'ia contaminada (Nota 5, Sec.~9).

\textbf{Las referencias tienen sentido porque est\'an anidadas.} El paseo aleatorio estacional es $(1-B)(1-B^{12})x_t=a_t$, es decir, el modelo de l\'ineas a\'ereas con $\theta=\Theta=0$. Su pron\'ostico es $f_{T,h}=x_T+(x_{T+h-12}-x_{T-12})$: ``el \'ultimo dato m\'as la variaci\'on del a\~no pasado''. El ingenuo estacional es $(1-B^{12})x_t=a_t$: ``igual que el mismo mes del a\~no pasado''. Ganarle a estas referencias muestra que estimar $\theta$ y $\Theta$ aporta.

\textbf{M\'etricas.}
\begin{itemize}
\item RMSE $=\sqrt{\tfrac1N\sum(Y-\hat Y)^2}$ y MAE: en unidades de la serie ($100\ln$, o sea \%).
\item MAPE: error porcentual sobre el \emph{nivel} (miles de viviendas). Se calcula en niveles y no sobre la serie centrada, porque dividir por valores cercanos a cero explota (Nota 5).
\item $U$ de Theil: entre 0 y 1, cuanto m\'as bajo mejor.
\end{itemize}

\textbf{Resultados.} El modelo de l\'ineas a\'ereas baja el RMSE del paseo aleatorio estacional un $\MejoraUno$, $\MejoraTres$, $\MejoraSeis$ y $\MejoraDoce\%$ (a 1, 3, 6 y 12 meses). A 12 meses empata con el ingenuo estacional. El MAPE a un mes es $\MapeFtresUno\%$.

\Q{¿Por qu\'e el MA(1) de la tasa no le gana al paseo aleatorio a 12 meses?}{Porque en un ARIMA(0,1,1) el pron\'ostico es $f_{t,1}=y_t-\theta a_t$ y despu\'es $f_{t,\ell}=f_{t,1}$ para todo $\ell\ge2$: es plano. La \'unica informaci\'on extra respecto del paseo aleatorio es la del primer paso. Es la funci\'on de pron\'ostico eventual de un $I(1)$ \nota{6}{3.4}.}

\Q{¿Para qu\'e el promedio de dos modelos?}{Es la combinaci\'on de pron\'osticos de Bates--Granger \nota{6}{11}. Combinar reduce el error si los errores de los modelos no est\'an muy correlacionados. Ac\'a F2 y F3 casi coinciden, as\'i que combinar no agrega nada. Es un resultado en s\'i mismo.}

%======================================================================
\section{La relaci\'on entre las series}

\subsection{La regresi\'on espuria \nota{6}{3.3}}
En niveles, $\ln$ inicios sobre $\ln$ tasa da pendiente \emph{positiva} ($\NivPend$, $t=\NivT$): m\'as tasa, m\'as inicios. Es absurdo econ\'omicamente y es el s\'intoma cl\'asico de Granger--Newbold: dos series $I(1)$ que se mueven juntas por sus tendencias producen una $t$ ``significativa'' sin relaci\'on real. La $t$ no sigue una $t$ de Student porque el residuo es $I(1)$.

\textbf{La regla $R^2>d_W$.} Sin dummies: $R^2=\NivR<d_W=\NivDW$, no dispara. Pero un $d_W$ de 0,15 implica $\hat\rho_{e,1}\approx1-d_W/2=0{,}93$: residuos casi con ra\'iz unitaria. Con dummies mensuales: $R^2=\DumR>d_W=\DumDW$, dispara. Con las dos series filtradas igual, la pendiente pasa a ser negativa, no significativa en el mismo mes, y $d_W=\FilDW$.

\Q{Si $R^2<d_W$ sin dummies, ¿por qu\'e dec\'is que es espuria?}{La regla es un s\'intoma suficiente, no una condici\'on necesaria. Lo que importa es que $d_W$ sea casi cero ($\hat\rho\approx0{,}93$), el signo absurdo, y que al filtrar las dos series la relaci\'on desaparezca en el mismo mes. Agregando las dummies, que no deber\'ian cambiar nada importante, la regla s\'i se cumple.}

\subsection{Correlaci\'on cruzada de innovaciones y test de Haugh}
\textbf{Idea.} No puedo correlacionar las series en niveles ni en diferencias, porque la autocorrelaci\'on propia de cada una genera correlaciones cruzadas espurias. Uso las \emph{innovaciones} $\hat a_t$ de cada modelo univariado, que son ruido blanco: si las series son independientes, $r_k\sim N(0,1/n)$ para todo $k$, con banda $\pm\BandaCC$.

\textbf{Haugh.} $S=n\sum_{k=-M}^{M}r_k^2\sim\chi^2_{2M+1}$ bajo independencia. Con $M=6$: $p\ \HaughSeisP$, no son independientes. Separo por direcci\'on, con $r_k=\operatorname{corr}(\hat a^{\text{tasa}}_t,\hat a^{\text{inicios}}_{t+k})$:
\begin{itemize}
\item $k\ge0$ (la tasa anticipa a los inicios): $p=\HaughDocePos$ en toda la muestra; $p=0{,}047$ en 1986--2007.
\item $k<0$ (los inicios anticipan a la tasa): $p=\HaughDoceNeg$.
\end{itemize}
\textbf{Lectura}: hay retroalimentaci\'on. Cuando la construcci\'on sorprende para arriba, la tasa sube despu\'es, porque la Fed y el mercado leen la actividad. Despu\'es de 2008 la direcci\'on tasa $\to$ inicios se pierde: la tasa la manejan las compras de la Fed.

\Q{¿Esto invalida que la tasa sea la causa?}{No invalida el mecanismo, pero s\'i el supuesto de exogeneidad estricta que pide la funci\'on de transferencia (``X ex\'ogena, sin retroalimentaci\'on'', segunda parte). Las salidas son: (i)~restringirse a 1986--2007, donde la direcci\'on tasa $\to$ inicios es la significativa; (ii)~usar como causa el bono a 10 a\~nos, que depende menos de la vivienda. Es la limitaci\'on principal del par, y la digo yo antes de que me la pregunten.}

\subsection{Rezagos distribuidos con el mismo filtro}
\textbf{El argumento.} Si la relaci\'on verdadera es $x_t=v(B)y_t+n_t$, como $F(B)$ y $v(B)$ son polinomios en $B$, conmutan:
\[
F(B)x_t=F(B)v(B)y_t+F(B)n_t=v(B)\,[F(B)y_t]+F(B)n_t .
\]
Filtrar las dos series con el mismo $F(B)=(1-B)(1-B^{12})$ deja intacta la respuesta $v(B)$ y vuelve estacionarias a las dos series. Por eso no es espuria. Estimo $v_0,\dots,v_{12}$ por MCO.

\textbf{Lectura.} Como las dos series est\'an en $100\ln$, los $v_j$ son elasticidades. La respuesta acumulada $\sum_{j\le k}v_j$ es $\ElasTres$ (ee $\ElasTresee$) a 3 meses y $\ElasDoce$ (ee $\ElasDoceee$) a 12 meses.

\textbf{De elasticidad a puntos porcentuales.} Una suba de 1\,pp desde $\TasaUltima\%$ es un cambio de $100\ln(7{,}67/6{,}67)=\PPequiv\%$ en la tasa. Efecto: $-0{,}54\times14{,}0\approx\PPDoce\%$ de inicios a un a\~no. El IC 95\% es $(-0{,}54\pm1{,}96\times0{,}33)\times14\approx[\PPDoceLo;\ \PPDoceHi]\%$: el signo es el esperado, pero no se puede descartar cero.

\Q{¿Por qu\'e los errores est\'andar son ``orientativos''?}{Porque el ruido filtrado $F(B)n_t$ no es ruido blanco: si $n_t$ era estacionario, al diferenciarlo queda sobrediferenciado, con estructura MA. Eso se ve en $d_W=\RDdw$, lejos de 2, que indica autocorrelaci\'on negativa. MCO sigue siendo insesgado, pero sus errores est\'andar est\'an mal. La estimaci\'on correcta es la funci\'on de transferencia, que modela $n_t$ como un ARMA (segunda parte).}

\begin{clave}En niveles la relaci\'on es espuria y con el signo equivocado. Bien filtrada aparece una elasticidad negativa de alrededor de $-0{,}5$ a un a\~no, imprecisa, y con retroalimentaci\'on que hay que tratar en la segunda parte.\end{clave}

%======================================================================
\section{El c\'odigo: qu\'e hace y decisiones t\'ecnicas}
\textbf{Software.} Todo est\'a hecho en R (software libre, regla 12). Las salidas ``tipo EViews'' son tablas que arm\'e con el formato de EViews para que se lean igual; \textbf{no las gener\'o EViews}. Hay que decirlo si se pregunta.

\begin{center}\small
\begin{tabularx}{\linewidth}{lX}
\toprule
Script & Qu\'e hace \\
\midrule
\texttt{R/funciones.R} & Herramientas comunes: datos, ADF (\texttt{adf}), correlograma con $Q$ (\texttt{correlograma}), estimaci\'on ARMA (\texttt{ajustar}), diagn\'ostico, rejilla de identificaci\'on, sobreajuste y el dibujo de las salidas. \\
\texttt{01\_datos.R} & Lee los CSV de FRED, pasa la tasa semanal a mensual y controla huecos. \\
\texttt{02\_exploracion.R} & Gr\'aficos con contexto, volatilidad por per\'iodo, rango--mediana, perfil de Box--Cox, tests $F$ de estacionalidad. \\
\texttt{03\_estacionariedad.R} & Correlogramas por filtro y todos los ADF. \\
\texttt{04\_inicios.R} & Filtros rivales, rejilla, modelo final, diagn\'ostico, estabilidad, serie desestacionalizada. \\
\texttt{05\_tasa.R} & Lo mismo para la tasa. \\
\texttt{06\_pronostico.R} & Evaluaci\'on con origen m\'ovil y pron\'ostico hacia adelante. \\
\texttt{07\_relacion.R} & Regresi\'on espuria, correlaci\'on cruzada, Haugh, rezagos distribuidos. \\
\bottomrule
\end{tabularx}
\end{center}

\textbf{Decisiones t\'ecnicas que me pueden preguntar:}
\begin{itemize}
\item \textbf{Estimaci\'on}: m\'axima verosimilitud exacta (\texttt{stats::arima}) sobre la serie ya filtrada $W_t$, sin constante en el doble filtro. Es lo mismo que estimar ARIMA$(0,1,1)(0,1,1)_{12}$ sobre el nivel.
\item \textbf{Errores est\'andar}: R usa el hessiano num\'erico. EViews por defecto usa el producto externo de gradientes (OPG), as\'i que los errores est\'andar podr\'ian diferir un poco en EViews. Los coeficientes no cambian.
\item \textbf{Criterios}: en las salidas, por observaci\'on, como EViews ($-2\ell/n+k\ln n/n$). En las comparaciones uso las diferencias de SBC totales; el ranking es el mismo.
\item \textbf{$p$-valores del ADF}: de MacKinnon, con el paquete \texttt{urca}. Son los mismos que reporta EViews.
\item \textbf{Signo del MA}: R y EViews usan $1+\theta B$; las Notas, $1-\theta B$. En el texto del informe todo est\'a con la convenci\'on de las Notas.
\item \textbf{Reproducibilidad}: \texttt{Rscript 00\_correr\_todo.R} rehace todo, y los n\'umeros del informe se escriben solos en \texttt{salidas/tablas/valores\_XX.tex}.
\end{itemize}

%======================================================================
\section{Hoja de n\'umeros para saber de memoria}\label{sec:numeros}
\begin{center}\small
\begin{tabularx}{\linewidth}{Xr}
\toprule
Muestra & 1985-01 a 2026-08, 500 meses \\
Volatilidad de la tasa: Volcker / 1985--2007 & $\VolVolcker$ / $\VolNormal$ pp \\
$\hat\lambda$ rango--mediana: inicios / tasa & $\LamInicios$ / $\LamTasa$ \\
$F$ estacionalidad / cambio de patr\'on (inicios) & $\EstFv$ / $\EstFcv$ \\
ADF nivel inicios ($\tau_\mu$) / tasa ($\tau_\mu$) & $\AdfVivNivTM$ / $\AdfTasaNivTM$ (VC $\AdfVCtm$) \\
ADF sobre la diferencia: inicios / tasa & $\AdfVivDif$ / $\AdfTasaDif$ \\
Modelo de los inicios: $\hat\theta$ / $\hat\Theta$ / $\hat\sigma$ & $\ThetaV$ / $\SThetaV$ / $\SigmaV\%$ \\
At\'ipico abril 2020 & $\AOV$ ($\AOVpct\%$) \\
Ljung--Box residuos de los inicios (12/24/36) & $\LBVdoce$ / $\LBVveinticuatro$ / $\LBVtreintayseis$ \\
Modelo de la tasa: $\hat\theta$ / $\hat\sigma$ & $\ThetaT$ / $\SigmaT\%$ \\
SBC: F3 le gana a F2 / F1 por & $\FiltroDSBCdos$ / $\FiltroDSBCuno$ puntos \\
Pron\'ostico: mejora sobre el paseo aleatorio estacional ($h$=1/3/6/12) & $\MejoraUno$/$\MejoraTres$/$\MejoraSeis$/$\MejoraDoce\%$ \\
Regresi\'on en niveles: pendiente / $R^2$ / $d_W$ & $\NivPend$ / $\NivR$ / $\NivDW$ \\
Haugh $M=6$ / solo $k\ge0$ / solo $k<0$ & $\HaughSeisP$ / $\HaughDocePos$ / $\HaughDoceNeg$ \\
Elasticidad a 12 meses (ee) & $\ElasDoce$ ($\ElasDoceee$) \\
Efecto de +1 pp a 12 meses (IC 95\%) & $\PPDoce\%$ ($[\PPDoceLo;\PPDoceHi]\%$) \\
\bottomrule
\end{tabularx}
\end{center}
\end{document}
